\documentclass[12pt,a4paper]{article}

% ==================================================
% 1. CONFIGURATION
% ==================================================

% Enable / disable title page
\newif\ifshowtitlepage
\showtitlepagetrue
% \showtitlepagefalse

% Document information
\newcommand{\worktitle}{Sprint 0 Retrospective}
\newcommand{\worksubtitle}{NAO Robot Exercise System}
\newcommand{\teamname}{NAOMI}
\newcommand{\subject}{NAO Robot Exercise System}

\newcommand{\university}{Slovak University of Technology in Bratislava}
\newcommand{\faculty}{Faculty of Informatics and Information Technologies}

\newcommand{\academicyear}{2026/2027}
\newcommand{\location}{Bratislava}
\newcommand{\documentyear}{2026/2027}

% Team leaders
\newcommand{\leaderone}{Ing. Jakub Perdek, PhD.}
\newcommand{\leadertwo}{doc. Ing. Fedor Lehocki, PhD.}

% Team members
\newcommand{\memberone}{Bc. Akhmadinurov Amal}
\newcommand{\membertwo}{Bc. Danilov Denis}
\newcommand{\memberthree}{Bc. Dmitriev Anton}
\newcommand{\memberfour}{Bc. Gainetdinov Ruslan}
\newcommand{\memberfive}{Bc. Kuzin Timofei}
\newcommand{\membersix}{Bc. Tikhoblazhenko Alena}
\newcommand{\memberseven}{Bc. Zimnitski Aliaksei}

% Sprint and meeting information -- CHECK BEFORE MERGE
\newcommand{\sprintname}{Sprint 0}
\newcommand{\sprintperiod}{15 October -- 25 October 2026}
\newcommand{\meetingdate}{Monday, 26 October 2026}
\newcommand{\meetingtime}{16:00}
\newcommand{\meetingplace}{NAO robot laboratory, FIIT STU}
\newcommand{\participants}{All team members: A. Akhmadinurov, D. Danilov, A. Dmitriev, R. Gainetdinov, T. Kuzin, A. Tikhoblazhenko, A. Zimnitski}
\newcommand{\reportauthor}{Bc. Gainetdinov Ruslan}


% ==================================================
% 2. PACKAGES
% ==================================================

\usepackage[
    a4paper,
    top=3cm,
    bottom=2.7cm,
    left=2.8cm,
    right=2.8cm,
    headheight=45pt,
    headsep=20pt
]{geometry}

\usepackage[T1]{fontenc}
\usepackage[utf8]{inputenc}
\usepackage[english]{babel}
\usepackage{lmodern}

\usepackage{graphicx}
\graphicspath{{../../assets/}{assets/}}
\usepackage[table]{xcolor}
\usepackage{fancyhdr}
\usepackage{titlesec}
\usepackage{setspace}
\usepackage{lastpage}
\usepackage{array}
\usepackage{float}
\usepackage{caption}
\usepackage{booktabs}
\usepackage{tabularx}
\usepackage{csquotes}

% No bibliography in progress reports, so biblatex is not loaded.

\usepackage[hidelinks]{hyperref}


% ==================================================
% 3. COLORS
% ==================================================

\definecolor{primary}{HTML}{000000}
\definecolor{secondary}{HTML}{000000}
\definecolor{muted}{HTML}{000000}
\definecolor{linecolor}{HTML}{000000}
\definecolor{naomiTableHeader}{HTML}{17445A}
\definecolor{naomiTableStripe}{HTML}{EDF7FA}
\definecolor{naomiTableRule}{HTML}{66BBD0}

\newcolumntype{L}[1]{>{\raggedright\arraybackslash}p{#1}}
\newcolumntype{Y}{>{\raggedright\arraybackslash}X}


% ==================================================
% 4. TYPOGRAPHY
% ==================================================

\setstretch{1.2}
\setlength{\parindent}{0pt}
\setlength{\parskip}{7pt}

\titleformat{\section}{\Large\bfseries\color{primary}}{\thesection}{0.8em}{}
\titleformat{\subsection}{\large\bfseries\color{primary}}{\thesubsection}{0.8em}{}
\titleformat{\subsubsection}{\normalsize\bfseries\color{primary}}{\thesubsubsection}{0.8em}{}

\titlespacing*{\section}{0pt}{20pt}{10pt}
\titlespacing*{\subsection}{0pt}{14pt}{7pt}
\titlespacing*{\subsubsection}{0pt}{12pt}{5pt}

\captionsetup{font=small, labelfont=bf, justification=centering}


% ==================================================
% 5. HEADER AND FOOTER
% ==================================================

\newcommand{\naomiheaderfooter}{%
    \fancyhf{}
    \fancyhead[L]{\includegraphics[width=4cm,height=1.15cm,keepaspectratio]{horizontal.png}}
    \fancyhead[R]{%
        \begin{tabular}{r}
            {\footnotesize\bfseries \worksubtitle}\\[3pt]
            {\scriptsize Academic Year \academicyear}
        \end{tabular}}
    \fancyfoot[L]{{\scriptsize \teamname\enspace|\enspace\subject}}
    \fancyfoot[R]{{\scriptsize \thepage\enspace/\enspace\pageref{LastPage}}}
    \renewcommand{\headrulewidth}{0.7pt}
    \renewcommand{\footrulewidth}{0.4pt}%
}

\pagestyle{fancy}
\naomiheaderfooter
\fancypagestyle{plain}{\naomiheaderfooter}


% ==================================================
% 6. TITLE PAGE
% ==================================================

\newcommand{\makemytitlepage}{
    \begin{titlepage}
        \thispagestyle{empty}
        \centering

        % University logo
        \vspace*{0.3cm}
        \includegraphics[width=6cm,height=5cm,keepaspectratio]{vertical.png}

        \vspace{0.5cm}
        {\large\bfseries \university\par}
        \vspace{0.2cm}
        {\normalsize \faculty\par}

        % Project and document title
        \vspace{1.7cm}
        {\small\MakeUppercase{\worksubtitle}\par}
        \vspace{0.5cm}
        \begin{minipage}{0.95\textwidth}
            \centering
            {\fontsize{25}{32}\selectfont\bfseries \worktitle\par}
        \end{minipage}

        % Sprint period and author
        \vspace{0.6cm}
        {\normalsize \sprintperiod\par}

        % Team information
        \vfill
        \begin{minipage}{0.9\textwidth}
            \centering
            \begin{minipage}[t]{0.46\linewidth}
                \centering
                {\normalsize\bfseries Team Leaders\par}
                \vspace{0.25cm}
                {\small
                \leaderone\par
                \leadertwo\par}
            \end{minipage}%
            \hfill
            \begin{minipage}[t]{0.48\linewidth}
                \centering
                {\normalsize\bfseries Team Members\par}
                \vspace{0.25cm}
                {\small
                \memberone\par
                \membertwo\par
                \memberthree\par
                \memberfour\par
                \memberfive\par
                \membersix\par
                \memberseven\par}
            \end{minipage}
        \end{minipage}

        % Bottom information
        \vfill
        {\small \location\enspace\textbullet\enspace\academicyear\par}
    \end{titlepage}
}


% ==================================================
% 7. IMAGES
% ==================================================

% \imageplaceholder{Text}{Caption} -- until the screenshot is added
\newcommand{\imageplaceholder}[2]{
    \begin{figure}[H]
        \centering
        \begingroup
        \setlength{\fboxsep}{0pt}
        \fbox{
            \begin{minipage}[c][5cm][c]{0.85\textwidth}
                \centering
                {\Large\bfseries IMAGE}
                \vspace{0.3cm}

                {\normalsize #1}
            \end{minipage}
        }
        \endgroup
        \caption{#2}
    \end{figure}
}

% \documentimage{path}{caption}{label}
\newcommand{\documentimage}[3]{
    \begin{figure}[H]
        \centering
        \includegraphics[width=0.85\textwidth,keepaspectratio]{#1}
        \caption{#2}
        \label{#3}
    \end{figure}
}


% ==================================================
% 8. TABLES
% ==================================================

\newenvironment{naomitable}[2]{%
    \begin{table}[H]
        \centering
        \caption{#1}
        \label{#2}
        \begingroup
        \small
        \renewcommand{\arraystretch}{1.4}
        \setlength{\tabcolsep}{6pt}
        \rowcolors{2}{naomiTableStripe}{white}
}{%
        \endgroup
    \end{table}
}


% ==================================================
% 9. DOCUMENT
% ==================================================

\begin{document}

\ifshowtitlepage
    \makemytitlepage
\fi


% ==================================================
% MEETING DETAILS
% ==================================================

\section{Meeting Details}

At the end of \sprintname{}, the \teamname{} team met to evaluate the sprint
in a Scrum retrospective. \sprintname{} was a preparatory sprint focused on
organizing the team and understanding the project.

\begin{naomitable}{Retrospective meeting}{tab:meeting}
    \begin{tabularx}{\textwidth}{L{4.5cm} Y}
        \rowcolor{naomiTableHeader}
        \textcolor{white}{\bfseries Item} &
        \textcolor{white}{\bfseries Details} \\
        Date and time & \meetingdate, \meetingtime \\
        Place & \meetingplace \\
        Retrospective for & \sprintname{} (\sprintperiod) \\
        Participants & \participants \\
        Prepared by & \reportauthor \\
        \arrayrulecolor{naomiTableRule}\bottomrule[1pt]
    \end{tabularx}
\end{naomitable}


% ==================================================
% SPRINT OUTCOME
% ==================================================

\section{Sprint Outcome}

The team completed 16 of 17 tasks of \sprintname{}. One task was not
completed: \textbf{SCRUM-28 Request and secure faculty server}.
The faculty assigned the server to the team, but connecting to it and
securing it took more time than expected, so the task was moved to Sprint~1.

\documentimage
    {board.png}
    {Board of \sprintname{} at sprint closure}
    {fig:board}

The main results of the sprint were the team's working environment
(Jira project, GitHub organization, fork of the existing repository,
documentation repository with automatic publishing), the first team
methodologies and the project description and goals.


% ==================================================
% TEAM MEETINGS
% ==================================================

\section{Team Meetings in the Sprint}

\textbf{13 October 2026 (Microsoft Teams) -- sprint planning.}
The team defined the first tasks: organizing the team and its tools,
analyzing the existing project to understand what needs to be done,
and preparing questions to clarify with the supervisor.

\textbf{19 October 2026 (Microsoft Teams).}
The team went through the topics discussed with the supervisor on
16 October 2026, turned them into concrete tasks and assigned at least one
task to each team member. Topics that were still unclear were recorded as
questions for the next supervisor meeting.


% ==================================================
% COURSE OF THE MEETING
% ==================================================

\section{Course of the Meeting}

The meeting took place in the laboratory with the NAO robot. It started
with a short summary of \sprintname{}: the team went through the completed
tasks and discussed the one task that was not completed (SCRUM-28) and the
reason for the delay.

The team agreed that from Sprint~1 all work will follow the methodologies
created in \sprintname{}. In particular, every task planned into a sprint
will be estimated in story points according to the backlog management
methodology, so that the team velocity can be tracked in the following
sprints.

In the second part of the meeting, the team tried to connect to the real
NAO robot with the existing solution. The connection did not work out of
the box: the services had missing dependencies, and the robot address and
port could not be configured without editing the Docker Compose files.
These problems were turned into Sprint~1 tasks (SCRUM-51, SCRUM-52 and
SCRUM-57).

% Optional, as in the supervisor's example: what each member said
% (what went well, what was a problem, what they would like to change).


% ==================================================
% IMPROVEMENTS
% ==================================================

\section{Changes for Sprint 1}

Based on the retrospective, the team introduced the following changes
from Sprint~1:

\begin{itemize}
    \item Every task in a sprint has an epic, an assignee, a story point
          estimate and a priority. Estimates follow the scale in the
          backlog management methodology.
    \item Changes to code and documents are merged only through pull
          requests with a description (Changes, Reason, Verification) and
          approvals from other team members.
    \item New tasks are created right after supervisor meetings, so that
          the backlog reflects the agreed scope.
\end{itemize}

\end{document}